\section{Where Liability Reaches, and Where It Stops}
\label{sec:experiment-closes}
Section~\ref{sec:military-responsibility} left the vendor in a gap between two bodies of law written for other actors. One of the two has since been rewritten for software, and what the rewriting reaches can be stated in a sentence: a duty to disclose evidence, and a presumption against the party withholding it, in force where the harm is physical and the defendant is a manufacturer. Those two conditions between them exclude the cases this book is built on, and this section says how.

The European Union's recast product liability directive, adopted in 2024 and applying to products put on the market from December 2026, makes a manufacturer, importer or distributor — the directive's term is economic operator — liable without proof of fault. It also carries the instrument built for opacity. A defendant can be ordered to disclose the evidence a claim needs, and where a product's technical complexity makes the claimant's proof excessively difficult a court may presume the defect and the causal link once the claimant shows each is likely \autocite[Arts.~8–10]{eu2024productliability}. A mechanism designed for a system nobody outside can inspect is law rather than a proposal. Two features fix how far it travels. The defendant is that economic operator, and not a state that deploys what it sold. And the damage recovered is the closed list §\ref{sec:military-responsibility} named: death or personal injury, including medically recognised damage to psychological health; damage to property; and the destruction or corruption of data not used professionally \autocite[Art.~6(1)]{eu2024productliability}. The first of those is wider than it sounds, and it is where this book's cases enter if they enter at all: a discriminatory score or an unappealable refusal is recoverable once it has made somebody ill, and not before. The wrong itself — the injury to a right, and the economic loss that follows it — is outside the list.

The proposal that would have reached those was withdrawn. The Commission's 2022 AI Liability Directive would have added a fault-based route, carrying the same presumption of causality, beyond defect and beyond the product regime's list of damage. Member states and industry could not agree on its terms; the Commission withdrew it in February 2025 and struck it from the legislative docket that October \autocite{europeancommission2022liability}. Outside the enacted directive's reach — a harm to a right, an economic loss, a state as operator — what remains is the gap Matthias described, distributed across designer, trainer, deployer and user and governed by whatever each national law says about a case none of them was written for. And a compensation regime of either kind answers one gap of the four Santoni de Sio and Mecacci distinguish, culpability. Moral accountability, public accountability, and the duty to prevent a harm before it occurs are untouched by both \autocite{santonidesio2021four}.

The 2017 electronic-personhood debate approached the same gap from the
other side, by asking whether the system itself could be made the
answerable party. The European Parliament floated a legal status for
the most autonomous systems \autocite{europeanparliament2017civil}, and
more than 150 researchers, roboticists and legal scholars objected that
personhood for a system with no consciousness would let a manufacturer
offload liability onto a legal fiction with no assets of its own
\autocite{robotics2018openletter}. Giving a bearer assets
(§\ref{sec:reserve-behind}) answers the second half of that objection
and not the first. Responsibility follows causal control: the builder
answers for how it formed the thing, the operator for what it demanded,
and the bearer for what it decided knowing what it knew, and no
arrangement of accounts moves a design decision away from the party that
made it.

Nothing in this is a research problem. Apportioning causal contribution across a training pipeline after a harm has no method, and until it has one the apportionment is drafted rather than measured. Two drafts exist. One was enacted and confined to defect; the other was negotiated, withdrawn, and left on the record with its terms intact. Neither was rejected as unworkable, and the second's terms are available to whoever drafts the next one.

Every entry in the research agenda and its appendix closes by naming two things, the access it requires and the authority. This is where the access exists. A duty to disclose evidence a party would rather withhold, backed by a presumption against that party when it does, is enforceable now — for defect claims, against economic operators, over physical harm. That is the reach, and the confinement is the finding. Beyond it a reviewer can be granted access, and publish what it finds, and still have produced nothing anybody must act on until somebody is answerable for it.
